\documentclass[10pt,a4paper]{article}

\usepackage[utf8]{inputenc}
\usepackage[T1]{fontenc}
\usepackage[margin=1.4cm,top=1.1cm,bottom=1.1cm]{geometry}
\usepackage{amsmath,amssymb}
\usepackage{enumitem}
\usepackage{array}
\usepackage{setspace}
\usepackage{booktabs}
\setstretch{1.0}

\setlength{\parindent}{0pt}
\setlength{\parskip}{0.15em}

\pagestyle{empty}

% ligne de réponse pleine largeur, espacée
\newcommand{\repligne}{\par\vspace{0.42cm}\noindent\dotfill}
\newcommand{\zonecalcul}[1]{\par\vspace{0.15cm}\noindent{\small\emph{Espace de calcul :}}\par\vspace{#1}}

\begin{document}

% ============= EN-TETE =============
\begin{center}
\textbf{\large Fiche d'exercices --- Séance 4 : des données propres, préparer avant d'apprendre}\\[0.05cm]
{\small Introduction au ML -- Dr. El Hadji Bassirou TOURÉ -- DMI $\cdot$ FST $\cdot$ UCAD \quad $\bullet$ \quad \textbf{Travail à la main, sans machine}}
\end{center}

\vspace{-0.15cm}
\hrule
\vspace{0.25cm}

{\small Cette fiche reprend, exercice par exercice, les calculs de la Séance 4 : taux de manquants, imputation par la médiane, encodage one-hot, standardisation (z-score), le modèle des $k$ plus proches voisins à la main, le réglage de $k$ par la paire de scores, puis la fuite de données et le Pipeline $f = m \circ \varphi$. Répondre directement sur la fiche.}

\vspace{0.2cm}

\fbox{\begin{minipage}{0.96\textwidth}
\small\textbf{Aides numériques (valables pour toute la fiche).} Quelques racines carrées utiles :
\begin{center}
\renewcommand{\arraystretch}{1.2}
\begin{tabular}{c|cccccccc}
$u$ & $1$ & $2$ & $2{,}5$ & $4{,}09$ & $8$ & $10$ & $13{,}70$ & $1{,}049$\\
\midrule
$\sqrt{u}$ & $1{,}00$ & $1{,}41$ & $1{,}58$ & $2{,}02$ & $2{,}83$ & $3{,}16$ & $3{,}70$ & $32{,}4$\\
\end{tabular}
\end{center}
Médiane d'un nombre \emph{pair} de valeurs $=$ moyenne des deux valeurs centrales (valeurs triées). Distance euclidienne en dimension $d$ : $d(x,x') = \sqrt{\sum_{j=1}^{d}(x_j - x'_j)^2}$.
\end{minipage}}

\vspace{0.45cm}

%=============================================================================
\textbf{\large Exercice 1 -- Taux de manquants}

\vspace{0.15cm}

Au poste de santé de Tivaouane, un registre de $400$ patients a été audité. Après remplacement des zéros \emph{impossibles} par NaN, on compte les valeurs manquantes par colonne :

\vspace{0.15cm}
{\small
\renewcommand{\arraystretch}{1.7}
\begin{tabular}{|l|c|c|c|c|c|}
\hline
colonne & glycémie & tension & insuline & IMC & fièvre\\
\hline
nombre de NaN & $8$ & $28$ & $180$ & $12$ & $0$\\
\hline
$m_j$ & & & & & \\
\hline
\end{tabular}
}

\vspace{0.2cm}
(a) Rappeler la formule du taux de manquants $m_j$, puis remplir la dernière ligne du tableau.
\zonecalcul{1.0cm}

\vspace{0.2cm}
(b) On envisage de \emph{jeter} toute ligne contenant au moins un trou. En supposant que les colonnes à trous ne se chevauchent pas, combien de patients perdrait-on au plus ? Combien en resterait-il ? Quelle proportion du jeu cela représente-t-il ?
\zonecalcul{1.2cm}
Patients perdus (au plus) : \makebox[2.5cm]{\dotfill} \quad restants : \makebox[2.5cm]{\dotfill} \quad proportion perdue : \makebox[2.5cm]{\dotfill}

\vspace{0.2cm}
(c) Pour chacune des cinq colonnes, proposer le traitement le plus raisonnable (imputation par la médiane / aucune action), et justifier en une ligne.
\repligne
\repligne
\repligne

\vspace{0.5cm}

%=============================================================================
\textbf{\large Exercice 2 -- Médiane et imputation}

\vspace{0.15cm}

Une mini-colonne d'insuline ($\mu$U/mL) de huit patients, après remplacement des zéros impossibles, contient :
\[
92 \quad \text{NaN} \quad 145 \quad 110 \quad \text{NaN} \quad 130 \quad 88 \quad 410
\]

(a) Trier les valeurs observées (sans les NaN), puis calculer la médiane.
\zonecalcul{1.1cm}
Valeurs triées : \dotfill

\vspace{0.15cm}
médiane $=$ \makebox[2.5cm]{\dotfill}

\vspace{0.2cm}
(b) Donner la colonne après imputation par la médiane (les valeurs observées ne bougent pas).

\vspace{0.15cm}
{\small
\renewcommand{\arraystretch}{1.7}
\begin{tabular}{|c|c|c|c|c|c|c|c|c|}
\hline
position & $1$ & $2$ & $3$ & $4$ & $5$ & $6$ & $7$ & $8$\\
\hline
valeur imputée & & & & & & & & \\
\hline
\end{tabular}
}

\vspace{0.2cm}
(c) Calculer la moyenne des valeurs observées. Combien de valeurs observées la moyenne dépasse-t-elle ? Conclure : pourquoi la médiane est-elle préférable ici ?
\zonecalcul{1.3cm}
moyenne $=$ \makebox[2.5cm]{\dotfill} \quad la moyenne dépasse \makebox[1.2cm]{\dotfill} valeurs sur $6$.
\repligne
\repligne

\vspace{0.5cm}

%=============================================================================
\textbf{\large Exercice 3 -- Encodage one-hot}

\vspace{0.15cm}

Une colonne « groupe sanguin » prend quatre modalités : \textbf{A}, \textbf{AB}, \textbf{B}, \textbf{O} (rangées dans cet ordre alphabétique).

\vspace{0.2cm}
(a) Combien de colonnes binaires le one-hot produit-il ? \makebox[2cm]{\dotfill}

\vspace{0.2cm}
(b) Encoder les trois patients suivants (une case « chaude » par ligne) :

\vspace{0.15cm}
{\small
\renewcommand{\arraystretch}{1.7}
\begin{tabular}{|l|c|c|c|c|}
\hline
groupe & A & AB & B & O\\
\hline
« B » & & & & \\
\hline
« O » & & & & \\
\hline
« AB » & & & & \\
\hline
\end{tabular}
}

\vspace{0.2cm}
(c) Pourquoi ne pas coder simplement $\text{A}=1,\ \text{AB}=2,\ \text{B}=3,\ \text{O}=4$ ? Expliquer en termes d'\emph{ordre} et de \emph{distance}.
\repligne
\repligne
\repligne

\vspace{0.2cm}
(d) Calculer la distance euclidienne entre les vecteurs one-hot de « A » et de « B ». Que vaudrait-elle pour deux modalités \emph{quelconques} distinctes ? Qu'en conclure sur la géométrie produite par le one-hot ?
\zonecalcul{1.0cm}
$d(\text{one-hot A},\ \text{one-hot B}) = $ \makebox[3cm]{\dotfill}
\repligne

\vspace{0.5cm}

%=============================================================================
\textbf{\large Exercice 4 -- Standardisation (z-score)}

\vspace{0.15cm}

Sur le jeu d'\emph{entraînement}, deux colonnes ont les statistiques suivantes : \quad hémoglobine $\mu = 12{,}5$, $\sigma = 1{,}8$ (g/dL) \qquad glycémie $\mu = 120$, $\sigma = 28$ (mg/dL).

\vspace{0.2cm}
(a) Rappeler la formule du z-score, puis calculer celui d'un patient à hémoglobine $15{,}2$ g/dL. Interpréter en une phrase.
\zonecalcul{1.1cm}
$z = $ \makebox[2.5cm]{\dotfill}
\repligne

\vspace{0.2cm}
(b) Un patient du jeu de \emph{test} a une hémoglobine de $9{,}8$ g/dL. Calculer son z-score. \textbf{Question clé :} faut-il recalculer $\mu$ et $\sigma$ en incluant ce patient ? Justifier.
\zonecalcul{1.1cm}
$z = $ \makebox[2.5cm]{\dotfill}
\repligne
\repligne

\vspace{0.2cm}
(c) Un troisième patient a une glycémie de $162$ mg/dL. Calculer son z-score, et le comparer à celui de la question (a). Que permet cette comparaison, alors que les deux colonnes ont des unités différentes ?
\zonecalcul{1.0cm}
$z = $ \makebox[2.5cm]{\dotfill}
\repligne

\vspace{0.5cm}

%=============================================================================
\textbf{\large Exercice 5 -- Les $k$ plus proches voisins, à la main}

\vspace{0.15cm}

\fbox{\begin{minipage}{0.96\textwidth}
\small Cinq patients d'entraînement sont décrits par deux caractéristiques \textbf{déjà standardisées} $(z_1, z_2)$ et une classe (P $=1$, N $=0$). On veut classer un nouveau patient $\star = (0{,}0,\ 1{,}0)$.
\begin{center}
\renewcommand{\arraystretch}{1.15}
\begin{tabular}{c|cc|c}
voisin & $z_1$ & $z_2$ & classe\\
\midrule
$A$ & $1{,}0$ & $2{,}0$ & $1$\\
$B$ & $-1{,}0$ & $1{,}0$ & $0$\\
$C$ & $0{,}5$ & $-0{,}5$ & $1$\\
$D$ & $3{,}0$ & $0{,}0$ & $1$\\
$E$ & $-2{,}0$ & $-1{,}0$ & $0$\\
\end{tabular}
\end{center}
\end{minipage}}

\vspace{0.2cm}
(a) Calculer la distance euclidienne de chaque voisin à $\star$, puis donner le rang (du plus proche au plus lointain).

\vspace{0.15cm}
{\small
\renewcommand{\arraystretch}{1.9}
\begin{tabular}{|c|c|c|c|}
\hline
voisin & calcul de la distance & distance & rang\\
\hline
$A$ & & & \\
\hline
$B$ & & & \\
\hline
$C$ & & & \\
\hline
$D$ & & & \\
\hline
$E$ & & & \\
\hline
\end{tabular}
}

\vspace{0.2cm}
(b) Donner la prédiction pour $k = 3$ : voisins retenus \makebox[3.5cm]{\dotfill} \quad classes \makebox[2.5cm]{\dotfill} \quad $f(\star) = $ \makebox[1.2cm]{\dotfill}

\vspace{0.15cm}
(c) Donner la prédiction pour $k = 5$ : classes \makebox[3.5cm]{\dotfill} \quad $f(\star) = $ \makebox[1.2cm]{\dotfill}

\vspace{0.2cm}
(d) Pourquoi choisit-on presque toujours $k$ \emph{impair} en classification binaire ?
\repligne
\repligne

\vspace{0.25cm}
(e) \textbf{L'échelle est une condition de validité.} Trois patients sont décrits par (âge, insuline) en unités \emph{brutes} : Awa $= (28,\ 100)$, Cheikh $= (31,\ 300)$, Fatou $= (65,\ 108)$. Sur le train, âge : $\mu = 35$, $\sigma = 10$ ; insuline : $\mu = 150$, $\sigma = 100$.

\vspace{0.15cm}
(e.1) En unités brutes, calculer $d(\text{Awa}, \text{Cheikh})$ et $d(\text{Awa}, \text{Fatou})$. Qui paraît la plus proche d'Awa ?
\zonecalcul{1.3cm}
$d_{\text{brut}}(\text{Awa},\text{Cheikh}) = $ \makebox[2.5cm]{\dotfill} \quad $d_{\text{brut}}(\text{Awa},\text{Fatou}) = $ \makebox[2.5cm]{\dotfill} \quad plus proche : \makebox[2cm]{\dotfill}

\vspace{0.15cm}
(e.2) Standardiser les trois patients (z-score de chaque coordonnée), puis recalculer les deux distances. Qui est maintenant la plus proche d'Awa ?
\zonecalcul{2.2cm}
$d_{\text{std}}(\text{Awa},\text{Cheikh}) = $ \makebox[2.5cm]{\dotfill} \quad $d_{\text{std}}(\text{Awa},\text{Fatou}) = $ \makebox[2.5cm]{\dotfill} \quad plus proche : \makebox[2cm]{\dotfill}

\vspace{0.15cm}
(e.3) L'ordre des voisins s'est-il inversé ? Expliquer en une phrase ce que la standardisation a corrigé.
\repligne
\repligne

\vspace{0.5cm}

%=============================================================================
\textbf{\large Exercice 6 -- Le réglage de $k$ : diagnostic}

\vspace{0.15cm}

Sur un jeu de dépistage, le kNN a été entraîné puis évalué pour trois valeurs de $k$. Baseline majoritaire : $0{,}66$. \textbf{Attention au sens de lecture : un \emph{petit} $k$ donne un modèle \emph{complexe}.}

\vspace{0.15cm}
{\small
\renewcommand{\arraystretch}{1.9}
\begin{tabular}{|c|c|c|c|c|}
\hline
$k$ & accuracy train & accuracy test & écart train $-$ test & diagnostic\\
\hline
$1$  & $1{,}00$ & $0{,}71$ & & \\
\hline
$15$ & $0{,}80$ & $0{,}78$ & & \\
\hline
$300$ & $0{,}67$ & $0{,}66$ & & \\
\hline
\end{tabular}
}

\vspace{0.2cm}
(a) Remplir la colonne des écarts.

\vspace{0.15cm}
(b) Remplir la colonne diagnostic (sous-ajusté / bien ajusté / sur-ajusté), puis justifier chaque cas en une phrase, en s'appuyant sur l'écart \emph{et} sur la baseline.

\vspace{0.2cm}
$k = 1$ : \dotfill
\repligne

\vspace{0.15cm}
$k = 15$ : \dotfill
\repligne

\vspace{0.15cm}
$k = 300$ : \dotfill
\repligne

\vspace{0.2cm}
(c) Valeur de $k$ à retenir : \makebox[1.5cm]{\dotfill} \quad sur la foi de quel chiffre ? \dotfill

\vspace{0.2cm}
(d) Le réglage $k = 1$ a la meilleure accuracy d'entraînement ($1{,}00$). Expliquer pourquoi ce chiffre ne doit jamais servir à choisir un modèle, et faire le lien avec la profondeur de l'arbre (Séance 3).
\repligne
\repligne
\repligne

\vspace{0.5cm}

%=============================================================================
\textbf{\large Exercice 7 -- Fuite de données et Pipeline}

\vspace{0.15cm}

(a) Pour chaque protocole, dire s'il y a \emph{fuite} et, le cas échéant, laquelle : \textbf{variable fuyarde} (une colonne contient la réponse) ou \textbf{prétraitement fuyard} ($\mu$, $\sigma$, médianes appris hors du train).

\vspace{0.15cm}
{\small
\renewcommand{\arraystretch}{1.8}
\begin{tabular}{|p{11.3cm}|c|}
\hline
protocole & fuite ? laquelle ?\\
\hline
(1) Prédire une réadmission avec la colonne « durée du séjour courant ». & \\
\hline
(2) Standardiser tout le jeu, \emph{puis} faire le split train/test. & \\
\hline
(3) Imputer le train avec les médianes du train, puis le test avec \emph{ces mêmes} médianes. & \\
\hline
(4) Faire le split, puis appeler \texttt{fit\_transform} du scaler sur le jeu de \emph{test}. & \\
\hline
\end{tabular}
}

\vspace{0.25cm}
(b) \textbf{La chaîne $\varphi$ à la main.} Un patient du \emph{test} traverse le Pipeline $f = m \circ \varphi$. Les paramètres ont été appris sur le \emph{train} : $\text{med}_{\text{insuline}} = 120$ ; glycémie $\mu = 120$, $\sigma = 28$ ; insuline $\mu = 150$, $\sigma = 100$ ; âge $\mu = 35$, $\sigma = 10$. Remplir le tableau étape par étape.

\vspace{0.15cm}
{\small
\renewcommand{\arraystretch}{2.1}
\begin{tabular}{|l|c|c|c|}
\hline
étape & glycémie & insuline & âge\\
\hline
brut & $162$ & NaN & $28$\\
\hline
après imputation & & & \\
\hline
après z-score & & & \\
\hline
\end{tabular}
}

\vspace{0.15cm}
Chaque nombre injecté vient-il du train ou du test ? Pourquoi est-ce essentiel ?
\repligne
\repligne

\vspace{0.25cm}
(c) \textbf{Remettre la chaîne dans l'ordre.} Les cinq étapes d'une étude propre, en désordre : \quad (1) score sur le test ; \quad (2) \texttt{train\_test\_split} ; \quad (3) remplacer les zéros impossibles par NaN ; \quad (4) \texttt{fit} du Pipeline sur le train ; \quad (5) construire le Pipeline (imputation $\to$ standardisation $\to$ modèle).

\vspace{0.2cm}
Ordre correct : \makebox[5cm]{\dotfill}

\vspace{0.2cm}
Où une inversion de $2$ et $4$ (\texttt{fit} avant le split) créerait-elle une fuite, et de quel type ?
\repligne
\repligne

\end{document}
